% ============================================================
%  H. Høffding (udg.), "Udvalgte Stykker af dansk filosofisk Litteratur"
%  = Mindesmærker af Danmarks Nationallitteratur, III. Kjøbenhavn:
%    Gyldendal 1910. 306 s.
%  TRANSCRIPTION (Danish) — Høffding's EDITORIAL MATTER for section IV,
%  "Den psykologiske Skole (Treschow, Sibbern, Howitz)".
%
%  SCOPE. This file transcribes Høffding's own words — the section heading,
%  his introductory note § 1 on Treschow, and the headings of the two
%  Treschow excerpts — NOT the excerpted Treschow text itself. The 1810
%  essay is transcribed complete, from the original printing, at
%  ../../treschow/enslige-ting/transcription.tex. The point of this file is
%  the *transmission document*: what Høffding said about Treschow and what
%  he chose to reprint. See RESUME-NOTES.md for the collation of his
%  selection against the complete essay.
%
%  Scan: ~/bibliotek/Høffding, Harald/mindesmaerkerafd03ande.pdf
%        (Internet Archive; 320 PDF pp.; clean OCR — preferred).
%        A second copy, udvalgte-stykker.pdf (HathiTrust/Wisconsin), has
%        badly jumbled OCR; use it only for cross-checking page images.
%  OFFSET (VERIFIED): renderer page = printed + 8.
%        Section IV opens on printed p. 99 (unnumbered section opener,
%        signature "7*" at foot) = render p. 107.
%        NB: pdftotext's form-feed indexing differs from pdftoppm's by ~3;
%        trust the renderer offset above when pulling page images.
%  TYPE: ANTIQUA throughout. Proper names in SMALL CAPS (Niels Treschow,
%        F. C. Sibbern). Danish low–high quotes „…". Em-dashes as printed.
% ============================================================
\documentclass[12pt,a4paper]{article}
\usepackage[utf8]{inputenc}
\usepackage[T1]{fontenc}
\usepackage{libertinus}
\usepackage{libertinust1math}
\usepackage[danish]{babel}
\usepackage[protrusion=true,expansion=true]{microtype}
\usepackage{setspace}
\usepackage[margin=1.2in]{geometry}
\usepackage{fancyhdr}
\usepackage{hyperref}

\hypersetup{
  pdftitle={Udvalgte Stykker af dansk filosofisk Litteratur — Treschow},
  pdfauthor={Harald Høffding},
  hidelinks
}

\pagestyle{fancy}
\fancyhf{}
\rhead{Høffding (udg.), \emph{Udvalgte Stykker} (1910) — Treschow}
\cfoot{\thepage}

\setstretch{1.3}

\title{\textbf{Den psykologiske Skole}\\[4pt]
  \large Høffdings Indledning til Treschow\\[6pt]
  \normalsize i \emph{Udvalgte Stykker af dansk filosofisk Litteratur}}
\author{Harald Høffding (udgiver)}
\date{\small Mindesmærker af Danmarks Nationallitteratur, III.\\
  Kjøbenhavn: Gyldendal, 1910}

\begin{document}
\maketitle
\thispagestyle{fancy}

\bigskip

% === TEXT BEGINS (printed p. 99 = render p. 107; section opener, folio suppressed) ===

% printed p. 99
\begin{center}
  {\large IV}\\[10pt]
  {\Large DEN PSYKOLOGISKE SKOLE (TRESCHOW,\\ SIBBERN, HOWITZ).}
\end{center}

\begin{center}\rule{0.10\linewidth}{0.4pt}\end{center}

\noindent
1.~Den Interesse for Erfaring og for de Problemer, der opstaa paa
Erfaringens Grund, som Holberg og Sneedorff havde været besjælede af, kom
igen til Orde hos en Række af Mænd, der i det nittende Aarhundrede dyrkede
Psykologien og søgte dels at bringe den i nøje Forbindelse med
Naturvidenskaben, dels at behandle videregaaende filosofiske Spørgsmaal ud
fra Erfaringssynspunkter. Den første i Rækken er \textsc{Niels Treschow}
(1751---1833), Forfatter af den første selvstændige danske Psykologi. I sin
Filosofi hævdede han --- i Modsætning til den tyske Filosofis Opgaaen i
almene og abstrakte Begreber --- det Individuelle og Konkrete som det
egentligt Reale („Gives der noget Begreb eller nogen Ide om enslige Ting?"
Videnskabernes Selskabs Skrifter 1807), og det var ham derfor naturligt at
antage, at Forskellene mellem Arterne havde udviklet sig historisk og
successivt. De „moralske" Betænkeligheder ved denne Antagelse har han paa en
smuk Maade afvist (Elementer til Historiens Filosofi 1811).

\bigskip
\noindent
{\large A. DET INDIVIDUELLES UUDTØMMELIGHED}

% [Her følger uddraget af "Gives der noget Begreb eller nogen Idee om enslige
%  Ting?" (1810), trykt s. 99--108. Uddraget gengiver originalens trykte
%  s. 225--230 og 235--241 (første linje); s. 231--234 og 241--254 er udeladt,
%  uden markering. Fuld tekst: ../../treschow/enslige-ting/transcription.tex.
%  Uddraget slutter: "...rigtigen er det alligevel at giøre Forskiel paa denne
%  abstrakte Forestilling om det Enslige og den concrete."]

\bigskip
\begin{center}\rule{0.22\linewidth}{0.4pt}\end{center}

% printed p. 108
\noindent
{\large B. UDVIKLINGSLÆRE}

% KILDE (verificeret): "Elementer til Historiens Philosophie" (1811).
% Høffding angiver ingen kilde ved uddraget selv; hans § 1 henviser til værket.
% Verificeret ved søgning i nb.no's fuldtekstindeks: "Thibetanerne" og
% "Evolutions" forekommer KUN i Elementer (1811) blandt Treschows værker, og
% s. 111--112 her er ord for ord identisk med Elementer, trykt s. 80
% (scan-s. 92; billedverificeret). NB: originalen er Fraktur og bruger "ø";
% Høffding har normaliseret ortografien. Citér originalen, ikke uddraget.

\noindent
Vi see gierne, at de Særsyner, som i Erfarenhed forekomme, forklares paa nogen
naturlig Maade. Vor Videlyst bliver derved tilfredsstillet, hvilken, naar
Almagt eller en anden overnaturlig Aarsag fremstilles, bliver afviist eller
stødt tilbage. I Fornuftens Barndom kiender man ei den Forskiel mellem det
Naturlige og Overnaturlige: begge blande sig i det mindre dannede Menneskes
Forestillingsart, og ere det samme. Det er maaskee mindre en virkelig
Vildfarelse end en ei tydelig tænkt Idee. Men, hvorledes det end dermed maa
forholde sig, saa ligger deri hele Grunden, hvorfor de Gamle
% printed p. 109
saavel i alle sædvanlige som usædvanlige Tilfælde henvendte deres Tanker til en
Guds eller forborgen Krafts Indflydelse; da vi derimod altid troe, at de
nærmeste Aarsager maae findes i noget Synligt og Bekiendt, men ansee de høiere
Aarsager alene for skiulte og oversandselige.

Hvad Verdens, Dyrs og Planters første Udspring angaaer, saa ere vi derfor
tilbøielige til at antage en ubegribelig Skabelse, endskiøndt vi tilstaae, at i
deres følgende Forandringer og Forplantelse ere Naturens egne Kræfter efter
engang bestemte Love fornemmelig virksomme. Fortsetter man de blot physiske
Undersøgelser endnu noget videre, saa synes man vel at blive et Slags Muelighed
vaer til at indsee, hvorledes de saakaldte livløse Stoffer og Legemer ved en
medfødt Kraft kan have dannet sig selv, men mistvivler om paa samme Maade at
begribe organiske Legemers Frembringelse. Heraf reiser sig det indtil de seneste
Tider med fast almindeligt Bifald optagne Evolutions-System, hvis sidste Grund
ere de mange Slags Spirer eller Kimer, der alene ved en umiddelbar Skabelse
syntes at kunne have faaet Tilværelse. At disse ved en raa Materies efterhaanden
sig selv uddannende Kraft, kunde være frembragt, denne Mening forekom de Fleste
ligesaa urimelig som farlig og nærgrændsende til udtrykkelig Gudsfornægtelse.
Hvor dybt, tænkte man desuden, nedværdiges ei derved den menneskelige Natur, ja
ethvert levende Væsen, uagtet de Gamle ei fandt det fornedrende, men ærbødigen
dyrkede Jorden som deres Moder. Vi see dog, at saadan er Naturens Gang. Om
Almagten giør man sig det Begreb, at den paa engang kan frembringe Alt, det
Største med det Mindste. Man forundrer sig derfor med Rette, naar man erfarer,
at Alting bliver til af noget Foregaaende, som, naar vi betragter det Hele,
altid hvad Formen angaaer er ringere end det Følgende. Frøet eller Kimen endog
til Mennesker og Dyr er i Begyndelsen intet uden en ganske formløs Vædske:
forfølge vi den til sit Udspring, naar Øjet ikke mere formaaer, med Tankerne;
saa taber den sig endelig i Stoffer, som vi ei
% printed p. 110
kan adskille fra de almindelige. Men hvor lidet er disse os bekiendte? hvor
herligt er det Støv, vi bruge som Sindbillede paa Ringhed, og af Uvidenhed
foragte? Er den fælles Materie vel andet end den synlige Gud, end en
Aabenbarelse af det høieste Væsen selv?\footnote{Ved Materien maa man her
forstaae ei det Phænomen, der er en Gienstand for Erfarenhed, men dens
oversandselige Substrat, det sidste Subject for samme, dens evige Væsen selv.}
Ligesom af hans uendelige Forstand Ideerne ved en evig Fødsel udvikle sig,
saaledes ligger Sæden, Grundformerne til alle Ting i Materiens Skiød indhyllede,
men blive først i Tiden, og følgelig kun efterhaanden, til anskuelige Billeder
af hine Mønstre, som de stedse nærmere stræbe at udtrykke.

De første Mennesker ere ei umiddelbar oprundne enten af Havet eller Jorden. Men
i begge ligger dog den middelbare Grund baade til vor Materie og Form. Af Jorden
ere vi dannede. Thi førend det Stive afsondrede sig fra det Flydende, var endnu
intet Levende til, allermindst den ypperligste Skabning blandt alle levende.
Naturens Kræfter ere fra Guddommens egen Kraft ei forskiellige, undtagen for
saavidt, at denne i sig selv betragtet er uendelig, men i hine aabenbarer sig
paa en endelig Maade, og i sin Uendelighed alene som en uophørlig tiltagende
Størrelse. I denne Mening kan man sige, at Materien har dannet sig selv,
efterhaanden gaaet frem fra de mindre fuldkomne til den mest fuldkomne
menneskelige Form. Man kan ligeledes med Sandhed sige, at Gud, efter at have
frembragt Materien selv, eller ved dens Egenskaber giort sit eget Væsen saa
anskueligt som det kunde blive, deraf igien har frembragt alle Ting. Disse
Talemaader, som de Fleste holde for himmelvidt forskiellige, betyde dog egentlig
det samme, og give kun tvende menneskelige Forestillingsarter tilkiende.

Paa en endnu mere middelbar Maade er Vandet alle Tings Oprindelse. De Gamle, der
af en let undskyldelig Utaalmodighed sprang de nærmeste og mellemste Aarsager
forbi, for at
% printed p. 111
opsvinge sig til de høieste, misforstode disse i sig selv ganske rigtige
Setninger om vort Udspring. Enten lode de Mennesker og Guder umiddelbar opstaa
af Chaos, som er den oprindelige Vædske, eller ogsaa meente de, at ved Solens,
Luftens, Vandets og Jordens forenede Kræfter avledes de fuldkomneste Dyr og
Mennesker selv af den sidste, ligesom Skimmel og Maddiker af forraadnede
Legemer. Den første Mening beroede paa en mythisk Tradition og en poetisk
Fremstillelse af samme. Digterens Phantasie ligner det Chaos, de besynge, og
indeholder paa samme Maade Sæden til de herligste Fornuftens Verker. Den
omfatter det Hele med et Øiekast, men mærker ikke, at hvad der i dette
Genieblik, ligesom i den guddommelige Forstand, paa engang er til, først
efterhaanden for Sandserne kan blive virkeligt. Den anden Hypothese er vel en
Frugt af Erfarenhed og Eftertanke, men en umoden. Thi, endskiøndt de nyeste
Erfaringer deri synes at stemme overeens med de allerældste, at nogle blandt de
ringeste Plante- og Dyrearter uden sædvanlig Forplantning af deres Lige, ja
endog ganske nye, paa denne Maade kan fremkomme, som tilforn ikke vare, saa have
vi dog af Naturens ordentlige Løb ikke mindste Grund til at formode det samme om
de større og mere uddannede Organismer i de Levendes tvende Riger. Saadanne
Kæmpeskridt, som en pludselig Overgang fra det laveste til det høieste Trin af
Uddannelse, ere ei mindre stridige mod Philosophiens Grundsetninger end mod
Erfarenhed og Analogie. Thi af Begrebet om Tid og Endelighed følger aabenbar det
Modsatte. Hvorledes var nogen Tid muelig uden Orden og en continuerlig Fremgang,
hvori der intet Spring kan være? Hvorledes kan man tænke sig det uendelige Væsen
udtrykt i endelige Former, uden formedelst en evig Annærmelse til hiint; eller
fatte Tiden som Evighedens Billede, med mindre den er uafbrudt og eensformig
overalt?

Fra Stoffer, der synes raae, fordi man deri endnu ingen Form bliver vaer, skeer
altsaa en umærkelig Overgang til mere bestemte Skikkelser, hvis ordentlige Følge
eller
% printed p. 112
Rekke man, som jeg ogsaa ovenfor har gjort, af den Aarsag blot i Almindelighed
kan angive. Ved skarp betegnende Grændser at adskille dem gaaer derfor ikke an,
fordi Naturens Gang i de Afvigelser eller Udartninger, hvorved de forrige
Slægter og Arter forvandles til ganske nye, ei blot er krybende, men saa langsom
tillige, at hverken egen Erfarenhed eller Historiens alt for lidet Omfang setter
os istand til at følge dens næsten ubemærkelige Spor. Efter Aartusindes Forløb
give de mindst tvetydige Levninger af Dyr og Mennesker, der kunne ansees for
Stammefædre til de nærværende, skiøndt temmelig forskiellige, dog Anledning til
Indvendinger og Tvivl. Jeg skal derfor her ingen anføre. Man vil i
Naturkyndiges Skrifter finde dem i Mængde beskrevne, men de fleste usikkre,
allesammen omtvistede. Hvad der i Henseende til disse kan være rigtigt eller
urigtigt kan man, som enhver indseer, \emph{a priori} ei afgiøre, men Sagen
selv, som man deraf vil bevise, er alligevel ikke tvivlsom.

En umiddelbar Skabelse, være sig af virkelige Dyr og Planter eller deres
nærmeste Kimer, er tomme Ord. Er derimod Spørgsmaalet om en middelbar eller
naturlig Frembringelse, saa er ingen anden muelig end den anførte. Det
menneskelige Kiøn har følgelig i sin nærværende Form og Skikkelse ei fra
Begyndelsen af været til, skiøndt enhver Historie, den hellige saavel som
verdslige, naturligviis maa gaae ud fra dette Punkt. De specifiske
Kiendemærker, hvorved Naturhistorien nu adskiller dette Væsen fra andre, har det
ei altid havt, men antaget efterhaanden. Jeg tør dristig fremsette denne Mening,
ei som en blot Hypothese, men som unægtelig Følge af en ikke tilfældig, men
nødvendig Naturorden, hvis Rekke man ei kan afbryde uden Urimelighed og
Modsigelse. Det vilde derimod være Forvovenhed, om jeg vilde paastaae, at man
med samme Vished kan afskildre den menneskelige Gestalt i enhver foregaaende
Tilstand, eller bestemme til hvilken Art af nu bekiendte Skabninger vi tilforn
have henhørt, eller af hvilken vi maaskee ere udsprungne. Jeg er overbeviist
% printed p. 113
om, at Mennesket udgjør en Art eller Slægt for sig selv, og troer ingenlunde, at
Aben, som Thibetanerne efter \emph{Pallas}'s Beretning paastaae, eller noget
andet Dyr, hvor høit de end i al Evighed kan stige, nogensinde vil blive
Menneske eller just det samme, om endog mere, end det, som vi nu ere; men vel er
jeg af den Mening, at vore Stammeforældre, hvis Afmindelse ingen Historie, intet
Archiv kan have giemt, have giennemløbet en Livets Bane, hvori hverken de høiere
Sindsgaver kunde komme til den Modenhed, som i det skiønne og aldeles uddannede
Legeme, vi nu bære, eller den blotte Sandselighed var forsynet med alle fornødne
Redskaber til saa mange Slags, saa fine og saa levende Fornemmelser. Naturens
Mangfoldighed, den uudtømmelige Rigdom af Ideer i Guddommens Skiød, deres
Bestemmelse til en evig og stedse fortskridende Virkelighed tillader os ei den
Tanke, at alle Væsener er intet uden det samme, men paa forskiellige Trin af den
Høide, som de maae bestige. Hvad betyder da Uligheden mellem Skabninger af en
Art, der uagtet denne individuelle Forskiel dog hver paa sin egen Maade kan være
lige fuldkomne? Hvorfor skulde saadan Mangfoldighed ei ogsaa i Arter have Sted?

Have vi maaskee Aarsag til at skamme os ved saadan Oprindelse? Dette Spørgsmaal
kan besvares ganske kort. Vi behøve ei at gaae saa langt tilbage i Tiden for at
finde stærkere Grunde til vor Ydmygelse. Hvor mange Aar er det vel siden, at
Enhver af os har staaet langt under Dyrenes, ja selv den ringeste Plantes Liv?
Er det ei vigtigere at tænke paa hvad vi nu ere og engang ganske vist skal blive
end hvad vi vare: som om den stolte \emph{Ceder} med Haan vilde see ned paa den
unge Spire, der oprinder ved dens Rødder, eller vi til det Foster, hvis Hierte
neppe har begyndt at slaae i Moders Liv. Overalt har alt Levende samme Rod:
guddommelig er vor Herkomst, evig og uendelig vor Bestemmelse.

\begin{center}\rule{0.22\linewidth}{0.4pt}\end{center}

\bigskip
% printed p. 114
% [§ 2 indleder F. C. Sibbern; § 3 formodentlig Howitz. Ikke transskriberet her.]

\end{document}
